\section{The Objective Function}

\subsection{The Claim}

Science finds short descriptions of large bodies of observations. The right
way to formalize ``short description'' is Kolmogorov complexity. The right
way to formalize ``finds'' is Bayesian inference with the Solomonoff prior.
Everything else in this paper is unpacking those three sentences.

The formal statement. Let $H$ be a hypothesis---a theory, a model, a
program---and let $D$ be observed data. The objective function for science is:
\begin{equation}
  \text{maximize} \quad \log P(H \mid D)
  = \log P(D \mid H) + \log P(H) + \text{const}
  \label{eq:bayes}
\end{equation}
by Bayes' rule, where $P(D \mid H)$ is the likelihood and $P(H)$ is the
prior. This is uncontroversial as a framework. The only question is: what
prior?

\subsection{The Prior Problem}

If the prior $P(H)$ can be chosen arbitrarily, Bayes' rule is a framework,
not a complete theory. Two scientists with different priors will reach
different posteriors from the same data. The standard objection to
Bayesianism is that this makes science subjective.

The Kolmogorov/Solomonoff answer: the prior is not arbitrary. There is a
unique, principled choice grounded in the theory of computation.

\subsection{The Solomonoff Prior}

The Kolmogorov complexity $\K(x)$ of a string $x$ is the length of the
shortest program (on a fixed universal Turing machine $U$) that outputs $x$
and halts:
\[
  \K(x) = \min\{ |p| : U(p) = x \}.
\]
Solomonoff~\cite{solomonoff1964} proposed assigning to each hypothesis $H$
the prior:
\[
  P(H) \propto 2^{-\K(H)}.
\]
Simpler hypotheses---those with shorter generators---receive exponentially
higher prior probability. This is Occam's Razor, derived from first
principles rather than assumed as a philosophical preference.

The Solomonoff prior is \emph{universal}: for any computable probability
distribution $P$ over strings, there exists a constant $c_P$ such that
$M(x) \geq c_P \cdot P(x)$ for all strings $x$, where $M$ is the Solomonoff
measure. No computable prior assigns more probability to any hypothesis than
$M$ does, up to a fixed factor. The prior dominates every computable
alternative.

\subsection{Minimum Description Length}

Because $\K$ is uncomputable (Section~\ref{sec:kolmogorov}), the practical
version is the Minimum Description Length (MDL) principle~\cite{rissanen1978}.
Instead of the shortest program, use the shortest description under a fixed
code. The objective becomes:
\begin{equation}
  \text{minimize} \quad L(H) + L(D \mid H)
  \label{eq:mdl}
\end{equation}
where $L(H)$ is the description length of the hypothesis and $L(D \mid H)$
is the description length of the data given the hypothesis (the negative
log-likelihood under a specific coding scheme).

The two terms trade off directly. $L(H)$ penalizes complexity: a theory that
is too elaborate spends too many bits describing itself. $L(D \mid H)$
penalizes poor fit: a theory that predicts the data badly leaves many residual
bits to transmit. The optimum is the theory that compresses the combined
description most. Science is the search for short two-part codes for the world.

\subsection{Occam's Razor, Derived}

A consequence worth stating explicitly: two theories that predict the data
equally well are not equally good. The simpler one has higher posterior. A
more complex theory must predict the data \emph{better} to compensate for its
higher complexity penalty. The scientist who adds epicycles to save a theory is
doing something formally wrong: increasing $\K(H)$ without a commensurate gain
in $\log P(D \mid H)$.

The objective function is \emph{objective} in both senses of the word: it is
a target (something to minimize) and it is observer-independent (grounded in
the structure of computation, not in any particular scientist's preferences).
The Solomonoff prior is invariant across observers up to a fixed constant
determined by the choice of universal Turing machine---not by whoever is doing
the science.